% 第 5 章
\bichapter{布线与布线后优化}{Routing and Postroute Optimization}
\label{chap:route}

本章介绍 IC Compiler II 路由器 Zroute 的布线能力，以及布线后优化功能。Zroute 面向多核硬件架构，可高效处理 45~nm 及以下先进工艺设计规则与可制造性设计（DFM）任务。

本章包括以下主题：
\begin{itemize}
  \item Zroute 简介（Introduction to Zroute）
  \item 基本 Zroute 流程（Basic Zroute Flow）
  \item 布线前提条件（Prerequisites for Routing）
  \item 定义 Via（Defining Vias）
  \item 插入 Via Ladder（Inserting Via Ladders）
  \item 检查可布线性（Checking Routability）
  \item 布线约束（Routing Constraints）
  \item 布线应用选项（Routing Application Options）
  \item 时钟网络布线（Routing Clock Nets）
  \item 关键网络布线（Routing Critical Nets）
  \item 次级电源/地 pin 布线（Routing Secondary Power and Ground Pins）
  \item 信号网络布线（Routing Signal Nets）
  \item 网络屏蔽（Shielding Nets）
  \item 执行布线后优化（Performing Postroute Optimization）
  \item 分析并修复信号电迁移违例（Analyzing and Fixing Signal Electromigration Violations）
  \item 执行 ECO 布线（Performing ECO Routing）
  \item 在 GUI 中布线（Routing Nets in the GUI）
  \item 清理已布线网络（Cleaning Up Routed Nets）
  \item 分析布线结果（Analyzing the Routing Results）
  \item 保存布线信息（Saving Route Information）
  \item 推导掩模颜色（Deriving Mask Colors）
  \item 插入与移除 Cut Metal 形状（Inserting and Removing Cut Metal Shapes）
\end{itemize}

% ============================================================
\section{Zroute 简介}
\subsection*{Introduction to Zroute}

Zroute 包含五个布线引擎：全局布线（global routing）、轨道分配（track assignment）、详细布线（detail routing）、ECO 布线与布线验证。全局布线、轨道分配与详细布线可通过任务专用命令或自动布线命令调用；ECO 布线与布线验证则使用任务专用命令。

Zroute 主要特性包括：
\begin{itemize}
  \item 全布线步骤（含全局布线、轨道分配、详细布线）支持多核多线程；
  \item 真实连通模型：矩形相触即视为电气连通，无需导线中心线重合；
  \item 动态迷宫网格：可 off-grid 连接 pin，同时保留网格路由器的速度优势；
  \item 多边形管理器：识别多边形，并理解 DRC 面向多边形；
  \item 详细布线阶段并发优化设计规则、天线规则、导线与 via；
  \item 详细布线阶段并发插入冗余 via；
  \item 全局/轨道/详细布线内建 soft rule 支持；
  \item 时序与串扰驱动的全局布线、轨道分配、详细布线与 ECO 布线；
  \item 智能设计规则处理（合并冗余违例、智能收敛）；
  \item 带层约束与非默认布线规则（NDR）的 net group 布线；
  \item 时钟布线；
  \item 布线验证；
  \item 以 soft rule 方式优化 DFM 与可制造性（DFY）；
  \item 支持先进设计规则，如多重图案化（multiple patterning）。
\end{itemize}

% ============================================================
\section{基本 Zroute 流程}
\subsection*{Basic Zroute Flow}

图~\ref{fig:basic-zroute-flow} 给出基本 Zroute 流程，包括时钟布线、信号布线、DFM 优化与布线验证。

\begin{figure}[htbp]
\centering
\fbox{\parbox{0.85\textwidth}{\centering\vspace{1.5cm}
\textit{（原书 Figure 63：Basic Zroute Flow）}\\[0.3em]
\small 请参阅原版 PDF 插图；可将截图放入 \texttt{figures/} 目录后用 \texttt{\textbackslash includegraphics} 替换。
\vspace{1.5cm}}}
\caption{基本 Zroute 流程}
\label{fig:basic-zroute-flow}
\end{figure}

% ============================================================
\section{布线前提条件}
\subsection*{Prerequisites for Routing}

运行 Zroute 前，block 与物理库须满足：
\begin{itemize}
  \item \textbf{库要求}：Zroute 从 technology file 获取全部设计规则信息，布线前须确保规则已在 technology file 中定义。详见 \textit{Synopsys Technology File and Routing Rules Reference Manual}。
  \item \textbf{Block 要求}：
  \begin{itemize}
    \item 电源/地网络已在设计规划后、布局前布线完成（见 \textit{IC Compiler II Design Planning User Guide}）；
    \item 已完成时钟树综合与优化（见第~\ref{chap:cts}~章）；
    \item 估计拥塞可接受；
    \item 估计时序可接受（约 0~ns slack）；
    \item 估计最大电容与 transition 无违例。
  \end{itemize}
\end{itemize}
后三项可通过 \cmd{check\_routability} 验证布局可布线性（见「检查可布线性」）。

% ============================================================
\section{定义 Via}
\subsection*{Defining Vias}

路由器支持以下 via 定义类型：
\begin{itemize}
  \item \textbf{Simple via 与 simple via array}：simple via 为单 cut via，由 cut 层及 cut/金属层矩形的高宽指定；simple via array 为多 cut 的 simple via。可用于：带 NDR 的时钟/信号布线、冗余 via 插入、带先进 via 规则的 PG 布线。
  \item \textbf{Custom via}：由任意 Manhattan 多边形集合构成的多 cut、异形 via；\textbf{仅}可用于冗余 via 插入。
\end{itemize}

Via 定义来源：
\begin{itemize}
  \item 设计库关联 technology file 的 \texttt{ContactCode} 段；
  \item LEF 文件中的 via rule \texttt{GENERATE} 语句（定义 simple via array）；
  \item 用户自定义 via（\cmd{create\_via\_def}）。
\end{itemize}

若希望 Zroute 将某 via 定义用于信号布线，须满足：\appopt{is\_default} 为 \cmd{true}，且 \appopt{is\_excluded\_for\_signal\_routing} 为 \cmd{false}。此外，Zroute 还会使用 technology file fat via table 或 \cmd{create\_routing\_rule} 所定义 NDR 中显式指定的非默认 via。

\subsection{从 LEF 读取 Via 定义}
\subsubsection*{Reading Via Definitions from a LEF File}

使用 \cmd{read\_tech\_lef} 读取 LEF 中 via rule \texttt{GENERATE} 语句。支持如下 LEF 语法：
\begin{lstlisting}
via_rule_name GENERATE [DEFAULT]
   LAYER lower_layer_name
      ENCLOSURE lower_overhang1 lower_overhang2
   LAYER upper_layer_name
      ENCLOSURE upper_overhang1 upper_overhang2
   LAYER cut_layer_name
      RECT llx lly urx ury
      SPACING x_spacing BY y_spacing
\end{lstlisting}
不支持 \texttt{WIDTH} 与 \texttt{RESISTANCE} 语句。

\subsection{创建 Via 定义}
\subsubsection*{Creating a Via Definition}

使用 \cmd{create\_via\_def} 创建 via 定义；创建后可在设计中任意处使用（类似 technology file 的 \texttt{ContactCode}）。用户 via 定义保存在设计库中，仅该设计可用。可定义 simple via、simple via array 与 custom via。

\subsubsection{定义 Simple Via}
\subsubsection*{Defining Simple Vias}

\begin{lstlisting}
create_via_def
   -cut_layer layer
   -cut_size {width height}
   -upper_enclosure {width height}
   -lower_enclosure {width height}
   [-min_rows number_of_rows]
   [-min_columns number_of_columns]
   [-min_cut_spacing distance]
   [-cut_pattern cut_pattern]
   [-is_default]
   [-force]
   via_def_name
\end{lstlisting}

无需指定 enclosure 层：工具根据 technology file 中与 via 层相邻的金属层自动确定。默认创建单 cut via；定义 via array 时指定 \opt{-min\_rows}/\opt{-min\_columns} 与 \opt{-min\_cut\_spacing}。默认 cut pattern 为满阵列；可用 \opt{-cut\_pattern} 修改。覆盖已有定义须用 \opt{-force}。

\begin{lstlisting}
icc2_shell> create_via_def design_via1_HV \
  -cut_layer VIA12 -cut_size {0.05 0.05} \
  -lower_enclosure {0.02 0.0} -upper_enclosure {0.0 0.2}

icc2_shell> create_via_def design_via12 -cut_pattern "01 10"
\end{lstlisting}

用 \cmd{report\_via\_defs} 报告用户 via 定义。

\subsubsection{定义 Custom Via}
\subsubsection*{Defining Custom Vias}

\begin{lstlisting}
create_via_def
   -shapes { {layer {coordinates} [mask_constraint]} ... }
   [-lower_mask_pattern alternating | uniform]
   [-upper_mask_pattern alternating | uniform]
   [-force]
   via_def_name
\end{lstlisting}

\opt{-shapes} 须为一个 via 层与两个金属 enclosure 层指定形状；每层可有多个形状。

\begin{lstlisting}
icc2_shell> create_via_def design_H_shape_via \
   -shapes { {VIA12 {0.035 -0.035} {0.100 0.035}}
             {VIA12 {-0.100 -0.035} {-0.035 0.035}}
             {METAL1 {-0.130 -0.035} {0.130 0.035}}
             {METAL2 {-0.100 -0.065} {-0.035 0.065}}
             {METAL2 {0.035 -0.065} {0.100 0.065}}
             {METAL2 {-0.100 -0.035} {0.100 0.035}} }
\end{lstlisting}

双图案化时，可对形状或 enclosure 层分配 mask constraint，二者不可同时使用：
\begin{itemize}
  \item 对形状：在 \opt{-shapes} 中用 \cmd{mask\_one}/\cmd{mask\_two}/\cmd{mask\_three}；
  \item 对 enclosure 层：用 \opt{-lower\_mask\_pattern}/\opt{-upper\_mask\_pattern}（\cmd{uniform} 或 \cmd{alternating}），创建 via 时再指定首个形状的 mask。
\end{itemize}

\begin{lstlisting}
icc2_shell> create_via_def VIA12 \
   -shapes { {M1 {-0.037 -0.010} {-0.017 0.01} mask_two}
             {M1 {0.017 -0.010} {0.037 0.01} mask_one}
             {VIA1 {-0.037 -0.01} {0.037 0.01} mask_two}
             {M2 {-0.044 -0.010} {0.044 0.010} mask_one} }
icc2_shell> create_via -net n1 -via_def VIA12 \
   -origin {100.100 200.320}
\end{lstlisting}

% ============================================================
\section{插入 Via Ladder}
\subsection*{Inserting Via Ladders}

Via ladder（亦称 via pillar）是自 pin 层向上堆叠至上层、供路由器连接的堆叠 via。Via ladder 可降低 via 电阻，改善性能与电迁移鲁棒性。

使用 via ladder 前须：
\begin{itemize}
  \item 在设计的工艺数据中定义 via ladder 规则（见「定义 Via Ladder 规则」）；
  \item 在全局布线\textbf{之前}插入 via ladder。
\end{itemize}

插入方法：
\begin{itemize}
  \item 布线前优化中的自动插入；
  \item 基于约束的插入：\cmd{insert\_via\_ladders}；
  \item 手动插入：\cmd{create\_via\_ladder}。
\end{itemize}

\begin{seeAlsoBox}
查询 Via Ladder；移除 Via Ladder；控制 Via Ladder 连接。
\end{seeAlsoBox}

\subsection{定义 Via Ladder 规则}
\subsubsection*{Defining Via Ladder Rules}

Via ladder 规则定义各级的行数与每行 cut 数，可在 technology file 中定义，或用 \cmd{create\_via\_rule} 创建。以下三种方式可创建等价规则：

\textbf{方式 1}：在 technology file 的 \texttt{ViaRule} 段中定义（含 \texttt{cutLayerNameTbl}、\texttt{numCutRowsTbl}、\texttt{numCutsPerRowTbl}、\texttt{upperMetalMinLengthTbl}、\texttt{cutXMinSpacingTbl}、\texttt{cutYMinSpacingTbl}、\texttt{maxNumStaggerTracksTbl}、\texttt{forHighPerformance}、\texttt{forElectromigration} 等）。详见 \textit{Synopsys Technology File and Routing Rules Reference Manual}。

\textbf{方式 2}：命令行详细规格 + 属性：
\begin{lstlisting}
icc2_shell> create_via_rule -name VL1 -cut_layer_names {VIA1 VIA2} \
   -cut_names {CUT1A CUT2A} -cut_rows {2 2} -cuts_per_row {2 2}
icc2_shell> set_attribute [get_via_rules VL1] \
   upper_metal_min_length_table {L1 L2}
icc2_shell> set_attribute [get_via_rules VL1] \
   for_high_performance false
icc2_shell> set_attribute [get_via_rules VL1] \
   for_electro_migration true
\end{lstlisting}

\textbf{方式 3}：先创建空规则再设属性（\appopt{cut\_layer\_name\_table\_size}、\appopt{cut\_layer\_name\_table}、\appopt{num\_cuts\_per\_row\_table} 等）。

无论创建方式如何，工具对 via rule 一视同仁。可用 \cmd{get\_via\_rules}、\cmd{remove\_via\_rules}、\cmd{report\_via\_rules}、\cmd{list\_attributes -application -class via\_rule}，以及 \cmd{write\_tech\_file} 保存。

\subsection{为电迁移 Via Ladder 生成规则}
\subsubsection*{Generating Via Ladder Rules for Electromigration Via Ladders}

可用 \cmd{generate\_via\_ladder\_template -config\_file} 根据 XML 配置生成规则与约束脚本，输出：
\begin{itemize}
  \item 模板脚本（默认 \texttt{auto\_gen\_via\_ladder\_template.tcl}，可用 \opt{-template\_file} 改名）：含 \cmd{create\_via\_rule} 与 \cmd{set\_attribute}；
  \item 关联脚本（默认 \texttt{auto\_gen\_via\_ladder\_association.tcl}，可用 \opt{-association\_file} 改名）：含 \cmd{set\_via\_ladder\_candidate}，并将 pin 的 \appopt{is\_em\_via\_ladder\_required} 设为 \cmd{true}。
\end{itemize}

配置文件基本结构：
\begin{lstlisting}
<EmRule>
   <Template name="rule_name">
      <Layer name="layer_name" row_number="cuts_per_row"
         [max_stagger_tracks="count"] ... />
   </Template>
   <Pin name="pin_name">
      <Template name="rule_name"/>
   </Pin>
</EmRule>
\end{lstlisting}

工具从 technology file 推导 cut 层、cut 名与各金属层最小长度。

\begin{noteBox}
运行 \cmd{generate\_via\_ladder\_template} 前须打开 block，以确保可访问工艺数据。
\end{noteBox}

\begin{lstlisting}
icc2_shell> generate_via_ladder_template -config_file vl_config
\end{lstlisting}

\subsection{为性能 Via Ladder 生成规则}
\subsubsection*{Generating Via Ladder Rules for Performance Via Ladders}

使用 \cmd{setup\_performance\_via\_ladder} 生成性能 via ladder 规则与关联。默认：
\begin{itemize}
  \item 生成至 M5 的规则（\opt{-max\_layer} 可改）；输出 XML（默认 \texttt{auto\_perf\_via\_ladder\_rule.xml}）与规则脚本（默认 \texttt{auto\_perf\_via\_ladder\_rule.tcl}）；也可用 \cmd{generate\_via\_rules\_for\_performance}；
  \item 将规则关联到无 \appopt{dont\_use} 的库单元 pin（\opt{-dont\_use} 可包含 dont\_use 单元；\opt{-lib\_pins} 可限定 pin）；输出关联脚本（默认 \texttt{auto\_perf\_via\_ladder\_association.tcl}）；也可用 \cmd{associate\_performance\_via\_ladder}；
  \item 运行所生成的脚本。
\end{itemize}

\subsection{基于约束的 Via Ladder 插入}
\subsubsection*{Constraint-Based Via Ladder Insertion}

步骤：
\begin{enumerate}
  \item 定义 via ladder 规则；
  \item 用 \cmd{set\_via\_ladder\_rules} / \cmd{set\_via\_ladder\_constraints} 定义约束；
  \item \cmd{insert\_via\_ladders} 插入；
  \item \cmd{verify\_via\_ladders} 验证。
\end{enumerate}

\subsubsection{定义 Via Ladder 约束}
\subsubsection*{Defining Via Ladder Constraints}

约束指定在哪些 pin 上插入何种模板；可定义全局与 instance 级约束，冲突时 instance 级优先。另可用 \cmd{set\_via\_ladder\_candidate} 为特定 pin 指定候选（布线前优化在时序关键路径上插入时使用）。用 \cmd{report\_via\_ladder\_constraints} 报告。

\textbf{全局约束}（via ladder rules）：\cmd{set\_via\_ladder\_rules}。映射选项：\opt{-master\_pin\_map}/\opt{-master\_pin\_map\_file}（格式 \texttt{\{\{lib\_cell/pin \{via\_ladder\_list\}\}...\}}）、\opt{-default\_ladders}。作用范围：\opt{-all\_instances\_of}、\opt{-all\_clock\_outputs}、\opt{-all\_clock\_inputs}、\opt{-all\_pins\_driving}。报告/移除：\cmd{report\_via\_ladder\_rules}、\cmd{remove\_via\_ladder\_rules}。

\textbf{Instance 级约束}：
\begin{lstlisting}
icc2_shell> set_via_ladder_constraints -pins {u1/i1 u2/i2} \
   {VL1 VL2 VL3}
\end{lstlisting}
移除用 \cmd{remove\_via\_ladder\_constraints}。

\subsubsection{插入 Via Ladder}
\subsubsection*{Inserting Via Ladders}

\cmd{insert\_via\_ladders} 默认行为与常用选项：
\begin{itemize}
  \item 默认不删除已有 ladder；用 \opt{-clean true} 先清除；
  \item 默认针对约束中全部 pin（信号与时钟）；用 \opt{-nets} 限定网络；
  \item 每个目标 pin 插入一个 ladder；多模板时选用首个不导致 DRC 的模板；全导致 DRC 时默认选 DRC 代价最低者；
  \item \opt{-ignore\_routing\_shape\_drcs true}：检查时忽略详细布线形状；
  \item \opt{-relax\_line\_end\_via\_enclosure\_rule true}：放松线端 cut enclosure；
  \item \opt{-relax\_pin\_layer\_metal\_spacing\_rules true}：ladder 完全包在 pin 内时放松 pin 层间距；
  \item \opt{-allow\_drcs false}：禁止产生 DRC 的插入；
  \item \opt{-strictly\_honor\_cut\_table true}：仅用模板中的 cut；
  \item \opt{-shift\_vias\_on\_transition\_layers true}：过渡层 cut 偏离 track 以改善 DRC；
  \item \opt{-ndr\_on\_top\_layer\_only true}：仅顶层遵循非默认宽度；
  \item \opt{-connect\_within\_metal}/\opt{-connect\_within\_metal\_for\_via\_ladder}：要求 pin 层上方 via 完全落在 pin（含延伸）内；
  \item \opt{-allow\_patching true}：允许 enclosure 打补丁以满足最小长度；
  \item \opt{-auto\_stagger true}：规则未定义最大 stagger track 时仍允许 stagger；
  \item \opt{-verbose true}/\opt{-user\_debug true}：报告插入状态与失败详情（内部/硬/软 DRC）。
\end{itemize}

\begin{noteBox}
若 pin 过小无法容纳任一指定 ladder，则不插入。
\end{noteBox}

\subsubsection{保护 Via Ladder}
\subsubsection*{Protecting Via Ladders}

\begin{lstlisting}
icc2_shell> set_app_options -name design.enable_via_ladder_protection \
   -value true
\end{lstlisting}
为 \cmd{true} 时，\cmd{remove\_shapes}、\cmd{remove\_vias}、\cmd{remove\_objects}、\cmd{set\_attribute}、\cmd{set\_via\_def}、\cmd{add\_to\_edit\_group} 仅处理不属于 via ladder 的对象；属于 ladder 的对象会警告并跳过。

\subsubsection{验证与更新}
\subsubsection*{Verifying and Updating Via Ladders}

\cmd{verify\_via\_ladders} 检查 ladder 是否匹配约束并正确连到 pin；可用 \opt{-nets} 限定。若插入时用了 \opt{-shift\_vias\_on\_transition\_layers true}，验证时也须使用，否则会误报。\opt{-report\_all\_via\_ladders false} 可将每类报告限制为 40 条。

更新方式：\cmd{refresh\_via\_ladders}；或在 \cmd{route\_group}/\cmd{route\_auto}/\cmd{route\_eco} 中由 \appopt{route.auto\_via\_ladder.update\_during\_route}（默认 \cmd{true}）自动更新。相关应用选项包括 \appopt{route.auto\_via\_ladder.connect\_within\_metal}、\appopt{route.auto\_via\_ladder.auto\_stagger}、\appopt{route.auto\_via\_ladder.report\_all\_via\_ladders}、\appopt{route.auto\_via\_ladder.update\_on\_route\_group\_nets\_only} 等。

\subsubsection{手动插入}
\subsubsection*{Manual Via Ladder Insertion}

\begin{enumerate}
  \item 用 \cmd{create\_shape}/\cmd{create\_via} 创建组成形状；
  \item 用 \cmd{create\_via\_ladder -shapes} 创建 ladder（名称形如 \texttt{VIA\_LADDER\_n}）。
\end{enumerate}
可选属性选项：\opt{-via\_rule}、\opt{-electromigration}、\opt{-high\_performance}、\opt{-pattern\_must\_join}、\opt{-pin}。

查询：\cmd{get\_via\_ladders}、\cmd{report\_via\_ladders}。移除：
\begin{lstlisting}
icc2_shell> remove_via_ladders *
\end{lstlisting}
\opt{-nets} 可限定网络。

% ============================================================
\section{检查可布线性}
\subsection*{Checking Routability}

布局完成后，用 \cmd{check\_routability} 检查 block 是否具备详细布线条件。默认检查：
\begin{itemize}
  \item \textbf{被阻塞的标准单元端口}：物理 pin 均不可达则视为 blocked。可达判定：pin 内有伸向邻层的 via、pin 层上存在不少于搜索范围的路径，或较短路径止于邻层 via。默认搜索范围为 2 倍 layer pitch（\opt{-standard\_cell\_search\_range}，最大 10）。\opt{-connect\_standard\_cells\_within\_pins true} 检查 via 是否完全在 pin 内。\opt{-check\_standard\_cell\_blocked\_ports false} 可关闭。
  \item \textbf{被阻塞的顶层/宏端口}：默认搜索距离为 10 倍 pitch（\opt{-blocked\_range}/\opt{-blocked\_range\_via\_side}，最大 40）。\opt{-obey\_direction\_preference true} 要求沿优选方向可达。\opt{-check\_non\_standard\_cell\_blocked\_ports false} 可关闭。
  \item \textbf{越界 pin}（\opt{-check\_out\_of\_boundary false} 可关）；
  \item \textbf{最小网格违例}（\opt{-check\_min\_grid false} 可关）；
  \item \textbf{非法 via 定义}（未着色 via array 行列不超过 20；custom/非对称 simple via 不超过 1000 cut；via 定义总数不超过 65535——不可关闭）；
  \item \textbf{非法 real metal via cut blockage}（\opt{-check\_via\_cut\_blockage false} 可关）；
  \item \textbf{最小宽度设置}（NDR 最小宽/屏蔽宽不超过 technology 最大宽——不可关闭）。
\end{itemize}

可选检查（需显式打开）：\opt{-check\_pg\_blocked\_ports}、\opt{-check\_redundant\_pg\_shapes}、\opt{-check\_routing\_track\_space}、\opt{-check\_frozen\_net\_blocked\_ports}、\opt{-check\_no\_net\_pins}、\opt{-check\_real\_metal\_blockage\_overlap\_pin}、\opt{-check\_lib\_via\_cut\_blockage}、\opt{-check\_shield}、\opt{-via\_ladder}。

Pin 连接控制：\opt{-obey\_access\_edges}、\opt{-access\_edge\_whole\_side}、\opt{-report\_no\_access\_edge}、\opt{-allow\_via\_rotation}。

Blocked port 检查会考虑 \cmd{set\_ignored\_layers}/\cmd{set\_routing\_rule}/\cmd{set\_clock\_routing\_rules} 的软/硬层约束及 \appopt{route.common.freeze\_layer\_by\_layer\_name} 等；\opt{-honor\_layer\_constraints false} 可忽略层约束。

错误数据默认名为 \texttt{check\_routability.err}（\opt{-error\_data} 可改），随 block 保存；可用 error browser 查看。

% ============================================================
\section{布线约束}
\subsection*{Routing Constraints}

布线约束在布线过程中提供指导。表~\ref{tab:routing-constraints} 概述各类约束。

\begin{table}[htbp]
\centering
\caption{布线约束类型}
\label{tab:routing-constraints}
\begin{tabularx}{\textwidth}{@{}l X@{}}
\toprule
\textbf{指导内容} & \textbf{说明} \\
\midrule
控制布线层 & \cmd{set\_ignored\_layers} 等，见「指定布线资源」 \\
更严格最小宽/间距 & 定义并应用 NDR（\cmd{create\_routing\_rule}） \\
控制穿越 voltage area & 电压域布线规则 \\
控制布线方向 & 层优选方向；区域用 preferred-direction-only / switch-preferred-direction routing guide \\
限制非优选方向边数 & maximum-number-of-pattern routing guide \\
控制 off-grid 布线 & 见「控制 Off-Grid 布线」 \\
禁止特定 net 布线 & \appopt{physical\_status} 设为 \cmd{locked} \\
禁止特定区域布线 & routing blockage \\
为顶层布线预留空间 & corridor routing blockage \\
限制布线到特定区域 & routing corridor \\
区域优先级 & access preference routing guide \\
控制布线密度 & utilization routing guide \\
鼓励 river routing & single-layer river routing guide \\
改善硬宏 pin 可布线性 & pin access routing guides \\
控制 block 边界周围布线 & perimeter constraint objects \\
Metal cut allowed/forbidden 规则 & boundary cell 插入时自动创建，或 \cmd{derive\_metal\_cut\_routing\_guides} \\
控制 pin 连接 & must-join、pin connection 类型、pin tapering \\
控制 via ladder 连接 & 见相关小节 \\
限制重布线范围 & rerouting mode \\
\bottomrule
\end{tabularx}
\end{table}

\subsection{定义 Routing Blockage}
\subsubsection*{Defining Routing Blockages}

Routing blockage 定义特定层上禁止布线的区域，Zroute 视其为硬约束。用 \cmd{create\_routing\_blockage} 创建，至少指定边界与层。

\begin{itemize}
  \item 矩形：\opt{-boundary \{ \{llx lly\} \{urx ury\} \}}；
  \item 直角多边形：\opt{-boundary \{ \{x1 y1\} \{x2 y2\} ... \}}，也可由 geo\_mask 等物理对象并集生成（多连通则生成多个 blockage）；
  \item 层：\opt{-layers}（金属、via 或 poly）。
\end{itemize}

默认名称 \texttt{RB\_objId}，阻止边界内所有网络；按 real metal 参与 RC 提取与 DRC，须满足最小间距。常用选项：\opt{-name}、\opt{-cell}、\opt{-zero\_spacing}（零间距且不作 real metal；防止 via 插入时须用此选项）、\opt{-net\_types}（\cmd{signal}/\cmd{clock}/\cmd{power}/\cmd{ground} 等）。

\begin{lstlisting}
icc2_shell> create_routing_blockage \
   -boundary { {0.0 0.0} {100.0 100.0} } \
   -net_types signal -layers M1
icc2_shell> create_routing_blockage \
   -boundary { {0.0 0.0} {100.0 100.0} } \
   -net_types signal -layers V1 -zero_spacing
icc2_shell> create_routing_blockage \
   -boundary { {0.0 0.0} {100.0 100.0} } \
   -net_types {power ground} -layers M2
\end{lstlisting}

\textbf{为顶层布线预留空间}：创建时加 \opt{-reserve\_for\_top\_level\_routing}。块级实现时作普通 blockage；frame view 提取时移除以允许顶层布线。

查询：\cmd{get\_routing\_blockages}、\cmd{get\_objects\_by\_location -classes routing\_blockage}。移除：\cmd{remove\_routing\_blockages}（\opt{-all} 移除全部 corridor）。

\subsection{定义 Routing Guide}
\subsubsection*{Defining Routing Guides}

Routing guide 为区域提供布线指令。用户定义的 guide 由 Zroute 遵从，但 Advanced Route 工具不遵从。

\begin{noteBox}
用户定义的 routing guide 由 Zroute 遵从；Advanced Route 工具不遵从。
\end{noteBox}

用 \cmd{create\_routing\_guide} 创建，须指定矩形边界 \opt{-boundary \{\{llx lly\} \{urx ury\}\}} 及用途相关选项。表~\ref{tab:user-routing-guides} 列出用户定义类型。

\begin{table}[htbp]
\centering
\caption{用户定义 Routing Guide}
\label{tab:user-routing-guides}
\begin{tabularx}{\textwidth}{@{}X l@{}}
\toprule
\textbf{用途} & \textbf{选项} \\
\midrule
仅在优选方向布线 & \opt{-preferred\_direction\_only} \\
切换优选方向 & \opt{-switch\_preferred\_direction} \\
限制特定布线模式出现次数 & \opt{-max\_patterns} \\
控制布线密度 & \opt{-horizontal\_track\_utilization} / \opt{-vertical\_track\_utilization} \\
区域优先级 & \opt{-access\_preference} \\
鼓励 river routing & \opt{-river\_routing} \\
\bottomrule
\end{tabularx}
\end{table}

默认名称 \texttt{RD\#n}；可用 \opt{-name}、\opt{-layers}、\opt{-cell}。

\subsubsection{控制布线方向}
\subsubsection*{Using Routing Guides to Control the Routing Direction}

\begin{lstlisting}
icc2_shell> create_routing_guide -boundary {{0.0 0.0} {100.0 100.0}} \
   -layers {M4} -preferred_direction_only
icc2_shell> create_routing_guide -boundary {{0.0 0.0} {100.0 100.0}} \
   -layers {M1 M2} -switch_preferred_direction
\end{lstlisting}

\begin{noteBox}
若 switch-preferred-direction 与 preferred-direction-only 重叠，前者优先。
\end{noteBox}

\subsubsection{限制非优选方向边数}
\subsubsection*{Using Routing Guides to Limit Edges in the Nonpreferred Direction}

\begin{lstlisting}
icc2_shell> create_routing_guide -boundary {{50 50} {200 200}} \
   -max_patterns {{M2 non_pref_dir_edge 2} {M4 non_pref_dir_edge 3}} \
   -layers {M2 M4}
\end{lstlisting}
当前支持的 pattern 为 \cmd{non\_pref\_dir\_edge}。超过阈值时 Zroute 报告违例。

\subsubsection{控制密度与优先级、River Routing}
\subsubsection*{Density, Access Preference, and River Routing}

\begin{lstlisting}
icc2_shell> create_routing_guide -boundary {{0.0 0.0} {100.0 100.0}} \
   -horizontal_track_utilization 50 -vertical_track_utilization 30
icc2_shell> create_routing_guide -boundary {{0 0} {300 300}} \
   -access_preference {M1 {wire_access_preference {{0 0} {2 1}} 0.2}
                          {via_access_preference {{0 0} {5 5}} 1.0}
                          {via_access_preference {{40 40} {45 45}} 0.5}}
icc2_shell> create_routing_guide -boundary {{0.0 0.0} {100.0 100.0}} \
   -river_routing -layers {M2}
\end{lstlisting}

强度 0--1；强度为 1 时视为硬约束，小于 1 的区域被忽略。用 \cmd{report\_routing\_guides} 报告。查询/移除：\cmd{get\_routing\_guides}、\cmd{remove\_routing\_guides -all}。

\subsection{推导 Routing Guide}
\subsubsection*{Deriving Routing Guides}

\subsubsection{Pin Access Routing Guide}
\subsubsection*{Deriving Pin Access Routing Guides}

在更小工艺节点复用硬宏时，可用 \cmd{derive\_pin\_access\_routing\_guides} 改善宏 pin 可达性：

\begin{lstlisting}
icc2_shell> derive_pin_access_routing_guides -cells myMacro \
   -layers {M2 M3} -x_width 0.6 -y_width 0.7
\end{lstlisting}

为指定宏在指定层创建 routing guide，并在宏周围创建带 pin 开孔的 metal blockage 及相邻 via 层 blockage。默认仅在优选方向开孔；\opt{-nonpreferred\_direction} 可在非优选方向也开孔；\opt{-pin\_extension} 可延伸开孔。命令会将宏 pin 的 \appopt{is\_rectangle\_only\_rule\_waived} 设为 \cmd{true}。

\begin{figure}[htbp]
\centering
\fbox{\parbox{0.85\textwidth}{\centering\vspace{1.2cm}\textit{（原书 Figure 64：Pin Access Routing Guides and Blockages）}\vspace{1.2cm}}}
\caption{Pin Access Routing Guide 与 Blockage}
\label{fig:pin-access-guides}
\end{figure}

移动宏后须重新生成；不可手动修改。

\subsubsection{Metal Cut Routing Guide}
\subsubsection*{Deriving Metal Cut Routing Guides}

若 technology file 定义了 metal cut allowed / forbidden preferred grid extension 规则且已有 boundary cell，可用：
\begin{lstlisting}
derive_metal_cut_routing_guides -add_metal_cut_allowed
\end{lstlisting}
\opt{-check\_only} 可检查；错误数据默认 \texttt{block\_name.err}（\opt{-error\_view} 可改）。

\subsection{控制 Block 边界周围布线}
\subsubsection*{Controlling Routing Around the Block Boundary}

用 \cmd{derive\_perimeter\_constraint\_objects} 沿边界放置金属形状、routing guide 与 blockage。边界变更后须重新生成。

\begin{itemize}
  \item \opt{-perimeter\_objects metal\_preferred}：沿优选方向边插入浮空金属（连续形状或用 \opt{-stub} 的 metal stub）；
  \item \opt{-short\_metal\_length}/\opt{-metal\_spacing}：垂直于非优选边的短金属；
  \item \opt{-perimeter\_objects mprg\_nonpreferred}：沿非优选边插入最大 pattern guide（阈值为 0）；
  \item 另可沿边界插入 routing blockage。
\end{itemize}

常用选项：\opt{-hierarchical}、\opt{-layers}、\opt{-spacing\_from\_boundary}、\opt{-width\_nonpreferred}、\opt{-spacing\_nonpreferred}、\opt{-check\_only}。

\begin{figure}[htbp]
\centering
\fbox{\parbox{0.85\textwidth}{\centering\vspace{1.0cm}\textit{（原书 Figure 65：Continuous Metal Shape Perimeter Constraint Objects）}\vspace{1.0cm}}}
\caption{连续金属形状边界约束对象}
\label{fig:perimeter-metal-cont}
\end{figure}

\begin{figure}[htbp]
\centering
\fbox{\parbox{0.85\textwidth}{\centering\vspace{1.0cm}\textit{（原书 Figure 66：Metal Stub Perimeter Constraint Objects）}\vspace{1.0cm}}}
\caption{Metal Stub 边界约束对象}
\label{fig:perimeter-metal-stub}
\end{figure}

\begin{figure}[htbp]
\centering
\fbox{\parbox{0.85\textwidth}{\centering\vspace{1.0cm}\textit{（原书 Figure 67：Metal Shape Perimeter Constraint Objects With Additional Short Shapes）}\vspace{1.0cm}}}
\caption{带短金属的边界约束对象}
\label{fig:perimeter-short-metal}
\end{figure}

\begin{figure}[htbp]
\centering
\fbox{\parbox{0.85\textwidth}{\centering\vspace{1.0cm}\textit{（原书 Figure 68：Routing Guide Perimeter Constraint Objects）}\vspace{1.0cm}}}
\caption{Routing Guide 边界约束对象}
\label{fig:perimeter-guide}
\end{figure}

\subsection{在特定区域内布线（Routing Corridor）}
\subsubsection*{Routing Nets Within a Specific Region}

Routing corridor 将布线限制在指定区域。全局布线视其为硬约束；轨道分配与详细布线视其为软约束（为修 DRC 可能略微越界）。

\begin{figure}[htbp]
\centering
\fbox{\parbox{0.85\textwidth}{\centering\vspace{1.0cm}\textit{（原书 Figure 69：Using a Routing Corridor）}\vspace{1.0cm}}}
\caption{使用 Routing Corridor}
\label{fig:routing-corridor}
\end{figure}

用 \cmd{create\_routing\_corridor} 创建，\cmd{add\_to\_routing\_corridor} 将 net 加入，\cmd{check\_routing\_corridors} 检查。

\subsection{使用非默认布线规则（NDR）}
\subsubsection*{Using Nondefault Routing Rules}

用 \cmd{create\_routing\_rule} 定义比 technology file 更严格的宽度、间距、via 间距、via 选择、mask、屏蔽等规则，再用 \cmd{set\_routing\_rule} 或 \cmd{set\_clock\_routing\_rules} 赋给网络。

\subsubsection{最小线宽与间距}
\subsubsection*{Defining Minimum Wire Width and Spacing Rules}

\begin{lstlisting}
create_routing_rule rule_name \
   -widths {layer1 width1 ...} \
   -spacings {layer1 spacing1 ...}
\end{lstlisting}

也可定义 soft spacing（带 weight）：
\begin{lstlisting}
-soft_spacings { {layer spacing weight} ... }
\end{lstlisting}

\begin{figure}[htbp]
\centering
\fbox{\parbox{0.85\textwidth}{\centering\vspace{1.0cm}\textit{（原书 Figure 70：Ignoring Nondefault Spacing Rule Violations）}\vspace{1.0cm}}}
\caption{忽略非默认间距规则违例}
\label{fig:ignore-ndr-spacing}
\end{figure}

为 soft spacing 分配修复努力：
\begin{lstlisting}
icc2_shell> set_app_options \
   -name route.common.soft_rule_weight_to_effort_level_map \
   -value { {low low} {medium medium} {high high} }
\end{lstlisting}
effort 取值：\cmd{off}/\cmd{low}/\cmd{medium}/\cmd{high}（\cmd{low} weight 不可用 medium/high effort）。

\subsubsection{最小 Via 间距与非默认 Via}
\subsubsection*{Via Spacing and Nondefault Vias}

\begin{lstlisting}
icc2_shell> create_routing_rule via_spacing_rule \
   -via_spacings {{V1 V1 2.3} {V1 V2 3.2}}
\end{lstlisting}

指定非默认 via：\opt{-cuts}（由工具按 technology 选择）或 \opt{-vias}（显式列出 via 类型、cut 数与方向 \cmd{R}/\cmd{NR}）。

\subsubsection{Mask 约束与屏蔽规则}
\subsubsection*{Mask Constraints and Shielding Rules}

\begin{lstlisting}
icc2_shell> create_routing_rule clock_mask1 \
   -mask_constraints {M4 mask_one M5 mask_one}
icc2_shell> create_routing_rule shield_rule \
   -shield_widths {M1 0.1 M2 0.1 M3 0.1 M4 0.1 M5 0.1 M6 0.3} \
   -shield_spacings {M1 0.1 M2 0.1 M3 0.1 M4 0.1 M5 0.1 M6 0.3}
\end{lstlisting}

\begin{noteBox}
为确保 DRC 收敛，双图案化 mask 约束应仅用于极少数时序关键网络。
\end{noteBox}

报告：\cmd{report\_routing\_rules -verbose}（\opt{-output} 可写出 Tcl）。移除：\cmd{remove\_routing\_rules}。修改须先删后重建。

\subsubsection{将 NDR 赋给网络}
\subsubsection*{Assigning Nondefault Routing Rules to Nets}

\begin{itemize}
  \item \cmd{set\_clock\_routing\_rules}：CTS 前赋给时钟网，CTS 会传播到新建时钟网；
  \item \cmd{set\_routing\_rule}：赋给信号网，或 CTS 后的时钟网。
\end{itemize}

表~\ref{tab:clock-ndr-restrict} 列出限制时钟规则分配的选项。

\begin{table}[htbp]
\centering
\caption{限制时钟布线规则分配}
\label{tab:clock-ndr-restrict}
\begin{tabularx}{\textwidth}{@{}X l@{}}
\toprule
\textbf{赋给} & \textbf{选项} \\
\midrule
特定时钟树 & \opt{-clocks} \\
连到 clock root 的网络 & \opt{-net\_type root} \\
连到 clock sink 的网络 & \opt{-net\_type sink} \\
时钟树内部网络 & \opt{-net\_type internal} \\
特定时钟网络 & \opt{-nets} \\
\bottomrule
\end{tabularx}
\end{table}

\begin{figure}[htbp]
\centering
\fbox{\parbox{0.85\textwidth}{\centering\vspace{1.0cm}\textit{（原书 Figure 71：Root, Internal, and Sink Clock Net Types）}\vspace{1.0cm}}}
\caption{Root / Internal / Sink 时钟网络类型}
\label{fig:clock-net-types}
\end{figure}

\begin{lstlisting}
icc2_shell> set_clock_routing_rules -rules NDR1 -net_type root
icc2_shell> set_clock_routing_rules -rules NDR2 -net_type internal
icc2_shell> set_clock_routing_rules -rules NDR3 -net_type sink
icc2_shell> set_clock_tree_options -root_ndr_fanout_limit 300
\end{lstlisting}

\begin{figure}[htbp]
\centering
\fbox{\parbox{0.85\textwidth}{\centering\vspace{1.0cm}\textit{（原书 Figure 72：Using a Fanout Limit for Selecting Root Nets）}\vspace{1.0cm}}}
\caption{用扇出限制选择 Root 网络}
\label{fig:root-fanout-limit}
\end{figure}

\begin{noteBox}
较小的 \opt{-root\_ndr\_fanout\_limit} 会增加使用 root NDR 的时钟网数量，可能加剧拥塞。
\end{noteBox}

多条 NDR 时优先级：\cmd{set\_routing\_rule} $>$ \cmd{set\_clock\_routing\_rules -nets} $>$ \cmd{set\_clock\_routing\_rules -clocks} $>$ 全局 \cmd{set\_clock\_routing\_rules}。

信号网示例：
\begin{lstlisting}
icc2_shell> set_routing_rule -rule WideMetal [get_nets CLK]
icc2_shell> set_routing_rule -default_rule [get_nets CLK]
\end{lstlisting}
另有 \opt{-no\_rule}、\opt{-clear}。报告：\cmd{report\_routing\_rules -of\_objects}、\cmd{report\_clock\_routing\_rules}。

\subsection{控制 Off-Grid 布线}
\subsubsection*{Controlling Off-Grid Routing}

默认仅当 technology 中 \texttt{onWireTrack}/\texttt{onGrid} 为 1 时才要求对齐 track。可用以下选项覆盖：
\begin{lstlisting}
icc2_shell> set_app_options \
   -name route.common.wire_on_grid_by_layer_name \
   -value {{M2 true} {M3 true}}
icc2_shell> set_app_options \
   -name route.common.via_on_grid_by_layer_name \
   -value {{V2 true}}
icc2_shell> set_app_options \
   -name route.common.extra_via_off_grid_cost_multiplier_by_layer_name \
   -value {{M2 0.5}}
\end{lstlisting}
有效代价 $=$ 基础代价 $\times (1+\mathrm{multiplier})$，multiplier 范围 0.0--20.0。

\subsection{Must-Join Pin 布线}
\subsubsection*{Routing Must-Join Pins}

Must-join pin 的各 terminal 作为同一网络连接。默认随机结构（图~\ref{fig:must-join-default}）；为改善电磁特性，可用 pattern-based must-join（经 via ladder 能力实现，图~\ref{fig:must-join-single}）。

\begin{figure}[htbp]
\centering
\fbox{\parbox{0.85\textwidth}{\centering\vspace{1.0cm}\textit{（原书 Figure 73：Default Must-Join Pin Connection）}\vspace{1.0cm}}}
\caption{默认 Must-Join 连接}
\label{fig:must-join-default}
\end{figure}

\begin{figure}[htbp]
\centering
\fbox{\parbox{0.85\textwidth}{\centering\vspace{1.0cm}\textit{（原书 Figure 74：Single-Level Pattern-Based Must-Join Pin Connection）}\vspace{1.0cm}}}
\caption{单层 Pattern-Based Must-Join}
\label{fig:must-join-single}
\end{figure}

多层结构（图~\ref{fig:must-join-multi}）：将 \appopt{route.auto\_via\_ladder.pattern\_must\_join\_over\_pin\_layer} 设为 2。当库 pin 有 \appopt{pattern\_must\_join} 且运行 \cmd{route\_auto}/\cmd{route\_group}/\cmd{route\_eco} 时自动使用。

\begin{figure}[htbp]
\centering
\fbox{\parbox{0.85\textwidth}{\centering\vspace{1.0cm}\textit{（原书 Figure 75：Multi-Level Pattern-Based Must-Join Pin Connection）}\vspace{1.0cm}}}
\caption{多层 Pattern-Based Must-Join}
\label{fig:must-join-multi}
\end{figure}

\begin{lstlisting}
icc2_shell> set_app_options -name \
   route.auto_via_ladder.pattern_must_join_max_number_stagger_tracks \
   -value {3 0}
\end{lstlisting}

\begin{figure}[htbp]
\centering
\fbox{\parbox{0.85\textwidth}{\centering\vspace{1.0cm}\textit{（原书 Figure 76：Pattern-Based Must-Join Pin Connection With Staggering）}\vspace{1.0cm}}}
\caption{带 Stagger 的 Pattern-Based Must-Join}
\label{fig:must-join-stagger}
\end{figure}

无法创建时报告 ``Needs pattern must join pin connection'' DRC。可用 \appopt{route.auto\_via\_ladder.update\_pattern\_must\_join\_during\_route} 控制是否在每次布线时更新。

\subsection{控制 Pin 连接与 Tapering}
\subsubsection*{Controlling Pin Connections and Pin Tapering}

\begin{lstlisting}
icc2_shell> set_app_options \
   -name route.common.connect_within_pins_by_layer_name \
   -value {{M1 via_all_pins}}
\end{lstlisting}

Mode：\cmd{off}（默认）、\cmd{via\_standard\_cell\_pins}、\cmd{via\_wire\_standard\_cell\_pins}、\cmd{via\_all\_pins}、\cmd{via\_wire\_all\_pins}。

\begin{figure}[htbp]
\centering
\fbox{\parbox{0.85\textwidth}{\centering\vspace{1.0cm}\textit{（原书 Figure 77--79：Pin Connection Restrictions）}\vspace{1.0cm}}}
\caption{Pin 连接限制示意（无限制 / via / via+wire）}
\label{fig:pin-connect-modes}
\end{figure}

Pin tapering：默认基于距离；也可用层基 tapering（\opt{-taper\_over\_pin\_layers}/\opt{-taper\_under\_pin\_layers} 等）。二者互斥，同时指定时采用层基设置。相关应用选项：\appopt{route.detail.pin\_taper\_mode}、\appopt{route.detail.use\_wide\_wire\_effort\_level}、\appopt{route.detail.use\_wide\_wire\_to\_input\_pin} 等。

Via ladder 连接：默认连到 ladder 最顶层。\appopt{route.detail.via\_ladder\_upper\_layer\_via\_connection\_mode}、\appopt{route.detail.allow\_default\_rule\_nets\_via\_ladder\_lower\_layer\_connection} 可调整。

\subsection{设置重布线模式}
\subsubsection*{Setting the Rerouting Mode}

将 net 的 \appopt{physical\_status} 设为 \cmd{locked}（禁止重布）、\cmd{minor\_change}（仅小改）或 \cmd{unrestricted}（默认）。

\begin{lstlisting}
icc2_shell> set_attribute -objects [get_nets net1] \
   -name physical_status -value locked
\end{lstlisting}

% --- 补遗：走廊细节、层约束、天线与冗余 via ---

\subsection{定义与使用 Routing Corridor（详述）}
\subsubsection*{Defining and Using Routing Corridors}

Routing corridor 将特定网络的全局布线限制在一组相连矩形所定义的区域内，并可对每个矩形指定最小/最大布线层。主要用于在信号布线前布关键网。若 routing guide 与 corridor 冲突，corridor 优先。

\begin{lstlisting}
icc2_shell> create_routing_corridor -name corridor_a \
   -boundary { {10 10} {20 35} } \
   -min_layer_name M2 -max_layer_name M4
icc2_shell> create_routing_corridor_shape -routing_corridor corridor_a \
   -boundary { {20 25} {40 35} } \
   -min_layer_name M2 -max_layer_name M4
icc2_shell> add_to_routing_corridor corridor_a [get_nets {n1 n2}]
icc2_shell> check_routing_corridors RC_0
\end{lstlisting}

也可按路径创建：\opt{-path} + \opt{-width}；端帽样式 \opt{-start\_endcap}/\opt{-end\_endcap}（\cmd{flush}/\cmd{full\_width}/\cmd{half\_width}）。一个 net/supernet 只能属于一个 corridor，且 corridor 须覆盖其全部 pin。

修改：\cmd{create\_routing\_corridor\_shape}、\cmd{remove\_routing\_corridor\_shapes}、\cmd{add\_to\_routing\_corridor}、\cmd{remove\_from\_routing\_corridor}。报告：\cmd{report\_routing\_corridors}（\opt{-output} 可写出 Tcl）、\cmd{get\_routing\_corridors -of\_objects}。移除：\cmd{remove\_routing\_corridors -all}。

\subsection{布线层资源与冻结层（章内相关）}
\subsubsection*{Routing Layer Resources and Freeze Layers}

全局/网络/时钟的最小最大布线层分别由 \cmd{set\_ignored\_layers}、\cmd{set\_routing\_rule} 与 \cmd{set\_clock\_routing\_rules} 的 \opt{-min\_routing\_layer}/\opt{-max\_routing\_layer} 设置；软/硬模式由 \appopt{route.common.global\_min\_layer\_mode}、\appopt{route.common.net\_min\_layer\_mode} 等控制。冻结层由 \appopt{route.common.freeze\_layer\_by\_layer\_name} 与 \appopt{route.common.freeze\_via\_to\_frozen\_layer\_by\_layer\_name} 指定。电压域穿越与层优选方向的详细设置见第~2~章相关主题。

\subsection{天线规则与冗余 Via（章内要点）}
\subsubsection*{Antenna Rules and Redundant Vias}

详细布线阶段并发处理天线规则，并在报告中给出天线违例与冗余 via 转换率。用 \cmd{check\_routes} 的 \opt{-antenna} 可启用/关闭天线检查。冗余 via 可在信号布线期间并发插入；详细布线结束报告含 ``Redundant via conversion report''（按层与 weight 的优化转换率）。更完整的天线修复与冗余 via 专项流程见原书交叉引用主题及第~6~章芯片收尾相关内容。

% ============================================================
\section{布线应用选项}
\subsection*{Routing Application Options}

工具提供影响全部三个布线引擎的选项，以及分别影响全局布线、轨道分配、详细布线的选项。详见 man 页：\texttt{route.common\_options}、\texttt{route.global\_options}、\texttt{route.track\_options}、\texttt{route.detail\_options}。

执行布线命令时，已设置（或工具代设）的选项会写入布线日志。要显示全部布线选项（含未设置者）：
\begin{lstlisting}
icc2_shell> set_app_options \
   -name route.common.verbose_level -value 1
\end{lstlisting}

% ============================================================
\section{时钟网络布线}
\subsection*{Routing Clock Nets}

在布线其余网络前先布时钟网，使用 \cmd{route\_group}：
\begin{lstlisting}
icc2_shell> route_group -all_clock_nets
\end{lstlisting}
也可用 \opt{-nets} 指定部分时钟网。该命令对时钟网依次执行全局布线、轨道分配与详细布线。

\begin{itemize}
  \item 普通时钟网使用 A-tree 布线以最小化各 pin 到 driver 的距离；
  \item 时钟 mesh 默认 Steiner；设 \appopt{route.common.clock\_topology} 为 \cmd{comb} 可改用 comb；
  \item 默认忽略已有全局布线；\opt{-reuse\_existing\_global\_route true} 可增量复用；
  \item 若已有详细布线则做增量详细布线；
  \item 默认最多 40 次 search-and-repair；\opt{-max\_detail\_route\_iterations} 可改。
\end{itemize}

\begin{noteBox}
若违例已全部修复或判定无法再修，Zroute 会在达到最大迭代前提前停止。
\end{noteBox}

% ============================================================
\section{关键网络布线}
\subsection*{Routing Critical Nets}

\begin{lstlisting}
icc2_shell> route_group -nets collection_of_critical_nets
icc2_shell> route_group -from_file file_name
\end{lstlisting}

若网络关联 routing corridor，则在该区域内布线。全局布线视 corridor 为硬约束；轨道分配与详细布线视之为软约束。

默认忽略已有全局布线，复用 dangling wire 与已有详细布线。\opt{-reuse\_existing\_global\_route true} 做增量全局布线（与 corridor 不兼容）；\opt{-utilize\_dangling\_wires false} 禁用 dangling 复用；\opt{-stop\_after\_global\_route true} 在全局布线后停止。

默认不修 soft DRC；可设：
\begin{lstlisting}
icc2_shell> set_app_options \
   -name route.common.post_group_route_fix_soft_violations \
   -value true
\end{lstlisting}

% ============================================================
\section{次级电源与地 Pin 布线}
\subsection*{Routing Secondary Power and Ground Pins}

步骤：
\begin{enumerate}
  \item 确认标准单元 frame view 中次级 PG pin 属性正确；
  \item 设置布线约束；
  \item 用 \cmd{route\_group} 布线。
\end{enumerate}

\subsection{验证属性}
\subsubsection*{Verifying the Secondary Power and Ground Pin Attributes}

必需属性：\appopt{is\_secondary\_pg}=\cmd{true}；\appopt{port\_type}=\cmd{power} 或 \cmd{ground}。

\begin{lstlisting}
icc2_shell> report_attributes -application \
   [get_lib_pins -of_objects reflib/cell/frame -all \
      -filter "name == attr_name"]
\end{lstlisting}

\subsection{设置约束并布线}
\subsubsection*{Setting Constraints and Routing}

可与信号布线类似设置层约束、单连接、\appopt{route.common.connect\_within\_pins\_by\_layer\_name}、\appopt{route.common.number\_of\_secondary\_pg\_pin\_connections}（簇内最大 pin 数，0 表示不限）、NDR 等。若要将次级 PG 与 tie-off 分离约束，设 \appopt{route.common.separate\_tie\_off\_from\_secondary\_pg} 为 \cmd{true}。

Topology ECO 努力由 \appopt{route.detail.topology\_eco\_effort\_level} 控制：\cmd{off}（默认）/\cmd{low}（1 次）/\cmd{medium}（4 次）/\cmd{high}（8 次）。

\begin{lstlisting}
icc2_shell> route_group -nets VDD1
icc2_shell> add_buffer {TOP/U1001/Z} {libA/BUFHVT} \
   -new_cell_names mynewHVT
icc2_shell> legalize_placement
icc2_shell> connect_pg_net -automatic
icc2_shell> route_group -nets {VDD2}
\end{lstlisting}

% ============================================================
\section{信号网络布线}
\subsection*{Routing Signal Nets}

信号布线前，全部时钟网须已布线且无违例。方法：
\begin{itemize}
  \item 独立命令：\cmd{route\_global}、\cmd{route\_track}、\cmd{route\_detail}（适合定制流程或大块分步）；每条命令开始读入、结束更新 block，\textbf{不}检查输入数据完整性；
  \item 自动布线：\cmd{route\_auto}（依次全局、轨道、详细；适合验证收敛、拥塞与 DRC QoR）。
\end{itemize}

信号布线可插入冗余 via（见「在信号网上插入冗余 Via」）。

\subsection{全局布线}
\subsubsection*{Global Routing}

布线前设置 \texttt{route.common\_options} 与 \texttt{route.global\_options}。默认非时序驱动。全局路由器将 block 划分为 global routing cell（默认宽度等于标准单元高度，与 row 对齐），按 blockage、pin、track 计算容量与需求，报告 overflow。

先进节点可开 via 密度建模：
\begin{lstlisting}
set_app_options -name route.global.via_cut_modeling -value true
\end{lstlisting}

两阶段：phase 0 初始布线；随后多轮 rip-up/reroute 减拥塞。Effort：\cmd{minimum}/\cmd{low}/\cmd{medium}（默认，最多 3 轮）/\cmd{high}（4）/\cmd{ultra}（6）。\appopt{route.global.force\_full\_effort} 可强制跑满轮次。

结束后库中保存 g-link/g-via（供后续步骤；轨道分配与详细布线后移除）与拥塞数据。默认仅硬拥塞；\appopt{route.global.export\_soft\_congestion\_maps}=\cmd{true} 可保存软拥塞。

设计规划阶段探索：\cmd{route\_global -floorplan true}；层次化可用 \opt{-virtual\_flat all\_routing} 或 \cmd{top\_and\_interface\_routing\_only}。

\subsubsection{时序驱动全局布线}
\subsubsection*{Timing-Driven Global Routing}

\begin{lstlisting}
set_app_options -name route.global.timing_driven -value true
\end{lstlisting}
流程：算延时 $\rightarrow$ phase0 $\rightarrow$ 两轮减拥塞 $\rightarrow$ 更新时序 $\rightarrow$ 识别新关键网 $\rightarrow$ 再多轮重布关键网。Effort 与 \appopt{route.global.advance\_node\_timing\_driven\_effort}（\cmd{low}/\cmd{medium}/\cmd{high}）共同决定重布轮数（见表~\ref{tab:tdgr-phases}）。\appopt{route.global.timing\_driven\_effort\_level} 默认 \cmd{high}，偏重时序 QoR。

\begin{table}[htbp]
\centering
\caption{时序驱动全局布线重布轮数（示意）}
\label{tab:tdgr-phases}
\begin{tabular}{@{}llll@{}}
\toprule
\cmd{-effort\_level} & advance=\cmd{low} & \cmd{medium} & \cmd{high}（默认） \\
\midrule
low & 2 & 2 & 2 \\
medium（默认） & 3 & 3 & 5 \\
high & 4 & 4 & 6 \\
ultra & 6 & 6 & 8 \\
\bottomrule
\end{tabular}
\end{table}

\subsubsection{串扰驱动与增量全局布线}
\subsubsection*{Crosstalk-Driven and Incremental Global Routing}

\begin{lstlisting}
set_app_options -name route.global.crosstalk_driven -value true
set_app_options -name time.si_enable_analysis -value true
# 增量：
route_global -reuse_existing_global_route true
\end{lstlisting}

\subsection{轨道分配}
\subsubsection*{Track Assignment}

须先完成全局布线。\cmd{route\_track} 为各全局布线分配实际 track，并重布重叠导线。完成后不再保留全局布线。

\begin{noteBox}
轨道分配用实际金属形状替换全局布线后，block 中不再包含全局布线。
\end{noteBox}

时序/串扰驱动：\appopt{route.track.timing\_driven}、\appopt{route.track.crosstalk\_driven}（及 \appopt{time.si\_enable\_analysis}）。

\subsection{详细布线}
\subsubsection*{Detail Routing}

须先完成全局布线与轨道分配。详细路由器按分区查找 DRC，rip-up/reroute 修复，并并发处理天线规则、优化 via 数与线长。

默认：整块布线（\opt{-coordinates} 可限区域）；首轮均匀分区、后续按违例分布调整；最多 40 轮（\opt{-max\_number\_iterations}）；非时序驱动。相关选项：\appopt{route.detail.optimize\_partition\_size\_for\_drc}、\appopt{route.detail.drc\_convergence\_effort\_level}、\appopt{route.detail.force\_max\_number\_iterations}、\appopt{route.detail.timing\_driven}、\appopt{route.common.rc\_driven\_setup\_effort\_level}、\appopt{route.detail.repair\_shorts\_over\_macros\_effort\_level}、\appopt{route.common.post\_detail\_route\_fix\_soft\_violations}。

增量详细布线须用 \opt{-incremental true}（否则从 iteration 0 重开）：
\begin{lstlisting}
icc2_shell> route_detail -incremental true
\end{lstlisting}

\begin{noteBox}
增量详细布线不修复 open net；须用 ECO 布线（\cmd{route\_eco}）。
\end{noteBox}

Early mode：\cmd{set\_qor\_strategy -mode early -stage pnr}。可用 \appopt{route.detail.save\_after\_iterations} 在指定轮次保存 block。\appopt{route.detail.report\_ignore\_drc} 可从汇总中排除特定 DRC。

\subsection{自动布线}
\subsubsection*{Routing Signal Nets by Using Automatic Routing}

\begin{lstlisting}
icc2_shell> route_auto
\end{lstlisting}

默认忽略已有全局布线，顺序执行全局、轨道、详细，步骤间不保存。选项：\opt{-reuse\_existing\_global\_route}、\opt{-stop\_after\_track\_assignment}、\opt{-max\_detail\_route\_iterations}、\opt{-save\_after\_global\_route}/\opt{-save\_after\_track\_assignment}/\opt{-save\_after\_detail\_route}、\opt{-save\_cell\_prefix}。

时序驱动：将 \appopt{route.global.timing\_driven}、\appopt{route.track.timing\_driven}、\appopt{route.detail.timing\_driven} 均设为 \cmd{true}。串扰驱动：另开 \appopt{route.global.crosstalk\_driven}、\appopt{route.track.crosstalk\_driven} 与 \appopt{time.si\_enable\_analysis}。

% ============================================================
\section{网络屏蔽}
\subsection*{Shielding Nets}

Zroute 按屏蔽规则中的宽度与间距生成屏蔽导线。除同层屏蔽外，还可对上一层/下一层做同轴屏蔽（coaxial shielding，图~\ref{fig:coaxial-shield}），隔离更好但更耗布线资源。

\begin{figure}[htbp]
\centering
\fbox{\parbox{0.85\textwidth}{\centering\vspace{1.2cm}\textit{（原书 Figure 80：Coaxial Shielding）}\vspace{1.2cm}}}
\caption{同轴屏蔽}
\label{fig:coaxial-shield}
\end{figure}

\begin{itemize}
  \item \textbf{布线前屏蔽}（preroute）：覆盖更好，但可能造成信号布线拥塞；常用于关键时钟；
  \item \textbf{布线后屏蔽}（postroute）：对可布线性影响极小，保护较弱。
\end{itemize}

\begin{figure}[htbp]
\centering
\fbox{\parbox{0.85\textwidth}{\centering\vspace{1.2cm}\textit{（原书 Figure 81：Zroute Shielding Flow）}\vspace{1.2cm}}}
\caption{Zroute 屏蔽流程}
\label{fig:shield-flow}
\end{figure}

\subsection{定义并分配屏蔽规则}
\subsubsection*{Defining and Assigning Shielding Rules}

用 \cmd{create\_routing\_rule} 的 \opt{-shield\_widths}/\opt{-shield\_spacings} 定义；时钟网用 \cmd{set\_clock\_routing\_rules}（仅 CTS 前），信号网用 \cmd{set\_routing\_rule}。建议仅对高频/关键时钟用屏蔽，低频时钟用双倍间距。

\subsection{布线前屏蔽}
\subsubsection*{Performing Preroute Shielding}

时钟树布线后、信号布线前，用 \cmd{create\_shields} 按规则加屏蔽。默认：对所有有屏蔽规则且未 frozen 的网；忽略短于 4 pitch 的导线（\appopt{route.common.min\_shield\_length\_by\_layer\_name} 可按层改）；屏蔽接到地网（\opt{-with\_ground} 或 \appopt{route.common.shielding\_nets}）；同层屏蔽。

同轴屏蔽：\opt{-coaxial\_above}/\opt{-coaxial\_below}，配合 \opt{-coaxial\_above\_skip\_tracks} 等或 \opt{-coaxial\_skip\_tracks\_on\_layers} 或用户间距选项（不可混用）。\opt{-align\_to\_shape\_end}、\opt{-preferred\_direction\_only} 等可进一步控制。

\begin{lstlisting}
icc2_shell> create_shields -nets $clock_nets \
   -coaxial_below true -coaxial_below_skip_tracks 0 \
   -with_ground VSS
\end{lstlisting}

命令报告基于网络与线长的平均屏蔽率。若与 \cmd{report\_shields} 略有差异，以 \cmd{report\_shields} 为准（通常 $<$3\%，最多约 5\%）。

\subsection{信号布线中的软屏蔽规则}
\subsubsection*{Soft Shielding Rules During Signal Routing}

有屏蔽规则但未预屏蔽的网，Zroute 在全局/轨道/详细全程预留屏蔽空间，详细布线结束报告 soft spacing (shielding) 违例；\textbf{并不}真正插入屏蔽，须事后再跑 \cmd{create\_shields}。仅预留同层屏蔽空间。若希望有 PG 可用时不预留，设 \appopt{route.common.allow\_pg\_as\_shield}=\cmd{true}。

\subsection{布线后屏蔽与增量屏蔽}
\subsubsection*{Postroute and Incremental Shielding}

布线后同样用 \cmd{create\_shields}。若用 \opt{-nets} 且无 NDR，则用 technology 默认宽间距。产生 DRC 时自动做最多 5 轮详细布线；仍有违例可 \cmd{route\_detail -incremental true}；若破坏 tie-off，改用 \cmd{route\_eco}。

增量屏蔽：将 \appopt{route.common.reshield\_modified\_nets} 设为 \cmd{unshield} 或 \cmd{reshield}（默认 \cmd{off}）。仅支持工具内修改的网络，DEF 外部修改不支持。

查询屏蔽形状：\cmd{get\_shapes -shield\_only}。报告：\cmd{report\_shields}（\opt{-per\_layer true}）。布线中启用屏蔽检查：\appopt{route.common.check\_shield}=\cmd{true}。

\begin{lstlisting}
# 示例：时钟预屏蔽 + 关键网后屏蔽
create_routing_rule shield_rule \
  -shield_widths {M1 0.1 M2 0.1 M3 0.1 M4 0.1 M5 0.1 M6 0.3} \
  -shield_spacings {M1 0.1 M2 0.1 M3 0.1 M4 0.1 M5 0.1 M6 0.3}
set_clock_routing_rules -clocks CLK -rules shield_rule
synthesize_clock_trees
route_group -all_clock_nets -reuse_existing_global_route true
create_shields -nets $clock_nets -with_ground VSS
set_routing_rule -rule shield_rule $critical_nets
route_auto
create_shields -nets $critical_nets -with_ground VSS
\end{lstlisting}

% ============================================================
\section{执行布线后优化}
\subsection*{Performing Postroute Optimization}

两类优化：
\begin{itemize}
  \item \textbf{逻辑优化}：改善时序/面积/功耗 QoR，修逻辑 DRC，并做合法化与 ECO 布线——\cmd{route\_opt}；
  \item \textbf{可布线性优化}：加大单元间距以修 pin access 引起的布线 DRC——\cmd{optimize\_routability}。
\end{itemize}
若两者都跑，先逻辑后可布线性。

\subsection{布线后逻辑优化}
\subsubsection*{Performing Postroute Logic Optimization}

前置：设计已完全布线并合法化，逻辑/布线 DRC 不过多。用 \cmd{check\_legality}、\cmd{check\_routes} 验证。推荐流程：
\begin{enumerate}
  \item \cmd{compute\_clock\_latency} 更新时钟延时；
  \item 运行两次 \cmd{route\_opt}（第一次合法化/ECO 可能影响时序，第二次再优化）。
\end{enumerate}

\cmd{route\_opt} 步骤：提取并更新时序（默认 PrimeTime 延时引擎，需 O-2018.06 及以后 Library Manager 生成的库；可用 \cmd{check\_consistency\_settings}）$\rightarrow$ 启用的优化 $\rightarrow$ 合法化 $\rightarrow$ ECO 布线。

默认优化 data path 的 setup/hold/面积/逻辑 DRC。额外优化见表~\ref{tab:route-opt-enable}。

\begin{table}[htbp]
\centering
\caption{启用 \cmd{route\_opt} 额外优化的应用选项}
\label{tab:route-opt-enable}
\begin{tabularx}{\textwidth}{@{}X X@{}}
\toprule
\textbf{优化} & \textbf{应用选项} \\
\midrule
并发时钟与数据（CCD）及时钟逻辑 DRC & \appopt{route\_opt.flow.enable\_ccd}=\cmd{true} \\
时钟逻辑 DRC（未开 CCD 时） & \appopt{route\_opt.flow.enable\_cto}=\cmd{true} \\
时钟树面积回收 & \appopt{route\_opt.flow.enable\_clock\_power\_recovery}=\cmd{area} \\
时钟树功耗回收 & 同上 $=\cmd{power}$ \\
功耗优化 & \appopt{route\_opt.flow.enable\_power}=\cmd{true} \\
机器学习路径优化 & \appopt{route\_opt.flow.enable\_ml\_opto}=\cmd{true} 且 PBA=\cmd{path} \\
\bottomrule
\end{tabularx}
\end{table}

\begin{noteBox}
\cmd{route\_opt} 仅遵从 \texttt{route\_opt} 应用选项；除 \appopt{opt.common.allow\_physical\_feedthrough} 外不遵从 \texttt{opt} 选项。
\end{noteBox}

PBA：\appopt{time.pba\_optimization\_mode}=\cmd{path}（推荐）或 \cmd{exhaustive}。ECO 阶段默认 5 轮详细布线（\appopt{route.detail.eco\_max\_number\_of\_iterations}）；\appopt{route\_opt.eco\_route.mode}=\cmd{track} 可改做轨道分配。

签核时序违例可用 \cmd{set\_route\_opt\_target\_endpoints}（\opt{-setup\_timing}/\opt{-hold\_timing}）指定目标后跑 \cmd{route\_opt}，再 \opt{-reset}。Size-only 模式：\appopt{route\_opt.flow.size\_only\_mode} 取 \cmd{true\_footprint}/\cmd{footprint}/\cmd{equal}/\cmd{equal\_or\_smaller}。

\subsection{使用 hyper\_route\_opt}
\subsubsection*{Performing Postroute Optimization Using hyper\_route\_opt}

可用单次 \cmd{hyper\_route\_opt} 减少迭代次数；设置与 \cmd{route\_opt} 相同。可选 Tcl 回调 \cmd{snps\_hyper\_route\_opt\_post\_eco} 控制 metal fill、PG augmentation 等。之后仍可用 \cmd{set\_route\_opt\_target\_endpoints}+\cmd{route\_opt} 收尾，或 \cmd{route\_detail -incremental} 修剩余布线 DRC。

\subsection{修复 Pin Access 引起的 DRC}
\subsubsection*{Fixing DRC Violations Caused by Pin Access Issues}

\begin{lstlisting}
optimize_routability
# 预览：
optimize_routability -check_drc_rules
# 完成后移除 keepout：
optimize_routability -remove_keepouts
\end{lstlisting}

分析邻接单元违例，在违例侧加 keepout（默认 site 宽，\opt{-keepout\_width} 可指定；\opt{-flip} 可尝试翻转）。须随后 \cmd{route\_eco} 或 \opt{-route}；\appopt{route.common.eco\_fix\_drc\_in\_changed\_area\_only}=\cmd{on} 可仅在变更区修 DRC。

% ============================================================
\section{分析并修复信号电迁移违例}
\subsection*{Analyzing and Fixing Signal Electromigration Violations}

信号电迁移由线宽缩小与高速开关导致电流密度升高引起，可造成短路或开路。流程：
\begin{enumerate}
  \item \cmd{read\_signal\_em\_constraints} 加载约束（ITF 或 ALF）；
  \item \cmd{read\_saif} 或 \cmd{set\_switching\_activity} 标注开关活动；
  \item \cmd{report\_signal\_em} 分析；
  \item \cmd{fix\_signal\_em} 修复。
\end{enumerate}

\begin{lstlisting}
open_block block1_routed
read_signal_em_constraints em.itf
read_saif block1.saif
report_signal_em -violated -verbose > block1.signal_em.rpt
fix_signal_em
\end{lstlisting}

ITF：\opt{-encrypted} 可读加密文件。ALF：\opt{-format ALF}。约束以查找表存入设计库。

分析计算 average / RMS / peak 电流。相关应用选项见表~\ref{tab:em-appopts}（如 \appopt{em.net\_delta\_temperature}、\appopt{em.net\_global\_rms\_relaxation\_factor}、\appopt{em.net\_min\_duty\_ratio}、\appopt{em.net\_use\_waveform\_duration}、\appopt{em.net\_violation\_rule\_types} 等）。

\begin{table}[htbp]
\centering
\caption{信号电迁移分析应用选项（节选）}
\label{tab:em-appopts}
\begin{tabularx}{\textwidth}{@{}l l X@{}}
\toprule
\textbf{选项} & \textbf{默认} & \textbf{说明} \\
\midrule
\appopt{em.net\_delta\_temperature} & 5.0 & RMS 查找用 $\Delta T$ \\
\appopt{em.net\_duration\_condition\_for\_peak} & 0.0 & peak 检查 duration \\
\appopt{em.net\_global\_rms\_relaxation\_factor} & 1.0 & RMS 全局松弛因子 \\
\appopt{em.net\_use\_waveform\_duration} & false & 是否用波形 duration \\
\bottomrule
\end{tabularx}
\end{table}

\cmd{fix\_signal\_em} 通过加宽金属的 NDR 并做 ECO 布线修复；每调用一轮，有残留须再跑。可用 \opt{-nets} 限定。

% ============================================================
\section{执行 ECO 布线}
\subsection*{Performing ECO Routing}

修改网络后须用 \cmd{route\_eco} 重连，顺序执行全局、轨道、详细布线。默认考虑时序与串扰（会提取并更新时序）；若不需要可关闭各引擎的 timing/crosstalk\_driven 选项以加速。

默认：处理全部 open net（\opt{-nets} 可限定）；忽略已有全局布线（\opt{-reuse\_existing\_global\_route true} 可增量）；复用 dangling wire（\opt{-utilize\_dangling\_wires false} 可关）；最多 40 轮详细布线；在整块修硬 DRC。

\begin{itemize}
  \item \opt{-open\_net\_driven true}：仅在 open net 邻域修 DRC；
  \item \opt{-reroute modified\_nets\_only} 或 \cmd{modified\_nets\_first\_then\_others}；
  \item \appopt{route.common.post\_eco\_route\_fix\_soft\_violations}=\cmd{true} 修 soft DRC。
\end{itemize}

报告变更网络时默认前 100 条；\opt{-max\_reported\_nets -1} 报告全部。

% ============================================================
\section{在 GUI 中布线}
\subsection*{Routing Nets in the GUI}

在活动 layout 视图用 Create Route 工具交互布线（Edit 工具栏按钮，或 \textbf{Create $>$ Route}）。默认忽略 routing blockage（可在 Mouse Tool Options 中改为遵从）；使用 technology 宽间距而忽略 NDR（同样可在面板中开启遵从 NDR）。

\begin{noteBox}
Create Route 工具不遵从 \cmd{create\_routing\_guide} 定义的 routing guide；也不使用 NDR 中的屏蔽宽/间距。
\end{noteBox}

修改已布线网络的 GUI 工具：Area Push、Spread Wires、Stretch Connected、Quick Connect。面板上的 Help 可打开 man 页。

% ============================================================
\section{清理已布线网络}
\subsection*{Cleaning Up Routed Nets}

\begin{lstlisting}
icc2_shell> remove_redundant_shapes
icc2_shell> remove_redundant_shapes -nets my_net \
   -remove_loop_shapes true
\end{lstlisting}

默认基于设计视图中已存 DRC 信息移除 dangling/floating 形状；\opt{-initial\_drc\_from\_input false} 则先跑 \cmd{check\_routes}。可用 \opt{-nets}/\opt{-layers}/\opt{-route\_types} 限定。不改拓扑、不碰 terminal；不处理 open net 或有 DRC 的网。\opt{-remove\_dangling\_shapes false}、\opt{-remove\_floating\_shapes false}、\opt{-remove\_loop\_shapes true}、\opt{-report\_changed\_nets true}。清理后应重新做 DRC 验证。

% ============================================================
\section{分析布线结果}
\subsection*{Analyzing the Routing Results}

\subsection{拥塞报告与拥塞图}
\subsubsection*{Generating a Congestion Report and Map}

\begin{lstlisting}
icc2_shell> report_congestion
\end{lstlisting}

默认用库中拥塞图；若无则跑 congestion-map-only 全局布线。选项：\opt{-rerun\_global\_router}、\opt{-boundary}、\opt{-layers}、\opt{-mode global\_route\_cell\_edge\_based}、\opt{-include\_soft\_congestion\_maps}。

GUI：\textbf{View $>$ Map $>$ Global Route Congestion}。图~\ref{fig:cong-map}、图~\ref{fig:cong-grc} 示意拥塞图与 GRC 上的 overflow（如 18/9 表示需求 18、供给 9）。

\begin{figure}[htbp]
\centering
\fbox{\parbox{0.85\textwidth}{\centering\vspace{1.0cm}\textit{（原书 Figure 82：Global Route Congestion Map）}\vspace{1.0cm}}}
\caption{全局布线拥塞图}
\label{fig:cong-map}
\end{figure}

\begin{figure}[htbp]
\centering
\fbox{\parbox{0.85\textwidth}{\centering\vspace{1.0cm}\textit{（原书 Figure 83：Global Route Overflow on Global Routing Cell）}\vspace{1.0cm}}}
\caption{Global Routing Cell 上的 Overflow}
\label{fig:cong-grc}
\end{figure}

拥塞计算模式：各层 overflow 之和（默认，忽略 underflow），或总 demand 减总 supply（含 underflow，结果更乐观）。

\subsection{用 Zroute 做设计规则检查}
\subsubsection*{Performing Design Rule Checking Using Zroute}

\begin{lstlisting}
icc2_shell> check_routes
\end{lstlisting}

默认检查布线 DRC、open net、天线与 voltage area 违例（排除 user/frozen/PG 网）。可用 \opt{-nets}、\opt{-check\_from\_user\_shapes}、\opt{-check\_from\_frozen\_shapes}；\opt{-drc}/\opt{-open\_net}/\opt{-antenna}/\opt{-voltage\_area} 可分别关闭。\opt{-coordinates} 做区域检查（与 \opt{-nets} 互斥；区域检查不做 open/天线/tie-to-rail）。

错误写入 \texttt{zroute.err}。可用 DRC query 命令或 error browser 分析。基于 check\_routes 结果的增量详细布线：
\begin{lstlisting}
icc2_shell> route_detail -incremental true \
   -initial_drc_from_input true
\end{lstlisting}

\subsection{签核 DRC、外部工具与 LVS}
\subsubsection*{Signoff DRC, External Tools, and LVS}

\begin{itemize}
  \item 签核：\cmd{signoff\_check\_drc} / \cmd{signoff\_fix\_drc}（需 IC Validator 许可证）；
  \item 外部 Calibre：\cmd{read\_drc\_error\_file} 导入扁平 Calibre 错误文件；
  \item LVS：\cmd{check\_lvs} 检查 short/open/floating（\opt{-checks}、\opt{-nets}、\opt{-max\_errors}、\opt{-treat\_terminal\_as\_voltage\_source}、\opt{-open\_reporting detailed}、\opt{-report\_floating\_pins} 等）。
\end{itemize}

\begin{figure}[htbp]
\centering
\fbox{\parbox{0.85\textwidth}{\centering\vspace{1.0cm}\textit{（原书 Figure 84：Layout With Power and Ground Terminals）}\vspace{1.0cm}}}
\caption{含电源/地 Terminal 的版图}
\label{fig:pg-terminals}
\end{figure}

\subsection{报告布线结果与线长}
\subsubsection*{Reporting Routing Results and Wire Length}

\cmd{report\_design -routing}：各层信号/PG 形状数与线长、HV 分布、via 统计与 double via 转换率等（\opt{-hierarchical} 可含整棵物理层次）。

\cmd{report\_wirelength}：按层与方向汇总线长（含 global/detail、Virtual、BestAvailable；支持斜线）。

\subsection{DRC 查询命令}
\subsubsection*{Using the DRC Query Commands}

\begin{lstlisting}
icc2_shell> get_drc_error_data -all
icc2_shell> open_drc_error_data zroute.err
icc2_shell> get_drc_error_types -error_data zroute.err
icc2_shell> get_drc_errors -error_data zroute.err \
   -filter {type_name == "Diff net spacing"}
\end{lstlisting}

属性示例：\appopt{type\_name}、\appopt{bbox}、\appopt{objects}、\appopt{shape}。完整列表：\cmd{list\_attributes -application -class drc\_error}。

% ============================================================
\section{保存布线信息}
\subsection*{Saving Route Information}

\begin{lstlisting}
icc2_shell> write_routes
\end{lstlisting}

生成可复现金属形状与 via（含属性）的 Tcl 脚本；\textbf{不}包含 routing blockage（后者用 \cmd{write\_floorplan}）。

% ============================================================
\section{推导掩模颜色}
\subsection*{Deriving Mask Colors}

布线命令通常会分配 mask 颜色。对未着色形状，用 \cmd{derive\_mask\_constraint} 从最近底层 track 推导（宽线占多 track 时取与下一空闲 track 相反颜色）。

\begin{lstlisting}
icc2_shell> derive_mask_constraint \
   [get_shapes -of_objects [get_nets net1]]
icc2_shell> derive_mask_constraint -derive_cut_mask \
   [get_vias -within {{600 600} {700 700}}]
icc2_shell> derive_mask_constraint -follow_pin_mask \
   [get_shapes -of_objects [get_nets net1]]
icc2_shell> derive_mask_constraint -follow_wider_object_mask \
   [get_shapes -of_objects [get_nets net1]]
\end{lstlisting}

\begin{seeAlsoBox}
第~\ref{chap:icc2}~章「多重图案化概念」。
\end{seeAlsoBox}

% ============================================================
\section{插入与移除 Cut Metal 形状}
\subsection*{Inserting and Removing Cut Metal Shapes}

Cut metal 用于双图案化设计以减小线端最小间距。Cut metal 层及尺寸在 technology file 的 Layer 段由 \texttt{cutMetalWidth}、\texttt{cutMetalHeight}/\texttt{cutMetalExtension} 定义；与同 mask 号金属层对应（如 cutMetal1 $\leftrightarrow$ metal1）。

布线定稿后：
\begin{lstlisting}
icc2_shell> create_cut_metals
icc2_shell> remove_cut_metals
\end{lstlisting}

插入条件：优选方向间距等于 \texttt{cutMetalWidth}、该处尚无 cut metal、金属层有对应 cut metal 层。可插在 net--net、net--PG、PG--obstruction 及子单元内金属之间，并传播颜色。

写出 DEF 须 \cmd{write\_def -version 5.8}；写出 GDSII/OASIS 须 \cmd{write\_gds}/\cmd{write\_oasis} 的 \opt{-connect\_below\_cut\_metal}。布线变更后须先 \cmd{remove\_cut\_metals} 再重新插入。
